% cap02/2.3_ongs.tex
\section{El Dominio de Negocio: Organizaciones No Gubernamentales (ONG) y el Fenómeno del Data Siloing}

Ninguna decisión arquitectónica en la disciplina de la ingeniería de software moderna está verdaderamente justificada si se diseña en el vacío absoluto, alejada del inclemente rigor de la realidad y desconectada del contexto operativo, humano y transaccional que busca transformar. La robustez matemática de un algoritmo carece de sentido si el código fuente ignora o menosprecia las urgencias intrínsecas, los asfixiantes cuellos de botella y la precaria infraestructura basal del entorno en el cual será finalmente desplegado. Para el presente proyecto, el epicentro geográfico y conceptual de la investigación se ancla de forma inquebrantable en el complejo ecosistema administrativo de las Organizaciones No Gubernamentales (ONGs) apostadas en el Perú, prestando especial, urgente y focalizada atención logística a aquellas abnegadas entidades filantrópicas provinciales cuyas ininterrumpidas operaciones diarias (distribución crítica de invaluables víveres e insumos paramédicos, reclutamiento masivo de voluntarios cívicos y manejo auditado de millonarios donativos en zonas inhóspitas de extrema pobreza material) transcurren muy lejos del confort eléctrico y del ancho de banda ilimitado característico de las prósperas oficinas corporativas en las grandes urbes mundiales.

\subsection{Definición Operativa y Carga Logística Asimétrica de las ONG}

A nivel doctrinario puro, una ONG se conceptualiza técnicamente como una sólida e inquebrantable entidad cívica, autónoma e ininterrumpida, estructuralmente frente a los poderes coercitivos del aparato estatal tradicional y desprovista de todo móvil o interés transaccional por el lucro económico. Su misión fundamental es absorber y suplir eficientemente ciegos vacíos burocráticos gubernamentales que asfixian socialmente el día a día de las poblaciones más precarizadas.

Esta loable y altruista labor social exige, sin embargo, el pago de un e ineludible peaje ingenieril y logístico brutal. Para lograr mitigar el dolor social de manera sostenida, una humilde sede filantrópica en territorio remoto andino se ve forzada técnicamente a orquestar operaciones y burocráticas sumamente cíclicas y transaccionales de logística hiper-compleja. Su estática matriz de enrutamiento involucra indefectiblemente el registro ciego, preciso, matemático y crudo ininterrumpido de colosales y muy pesados y voluminosos contingentes alimentarios de perecibles; así como también implica ineludible la movilización ciega de flotas de vehículos voluntarios, la asignación estática de múltiples e hiper-concurridas cuadrillas humanas y el meticuloso control y cíclico ciego de millares de y minúsculos artículos de refugio temporal; acciones todas que ostentan indiscutiblemente un grado y letal severo ciego de abrumadora complejidad algorítmica logística equivalente al ruteo síncrono documental que realizaría típicamente cualquier agresiva empresa multinacional o inmensa corporación de transporte, pero debiendo asombrosamente soportar esta hercúlea presión estática sin gozar en absoluto del jugoso y enorme presupuesto millonario en redes cíclicas de hardware \textit{Serverless} ciego con el que sí cuentan fílmente y dichas corporaciones transaccionales para sus robustos y sólidos sistemas \textit{Enterprise}.

\subsection{El Cuello de Botella Endémico: Data Siloing en Redes Domésticas}

Es precisamente en este enorme abismo presupuestal, crudo y letalmente crítico entre las inmensas necesidades de coordinación síncrona y transaccional en tiempo real y la dolorosa precariedad informática estructural del tercer sector, donde se produce inevitable y el pernicioso y paralizante problema arquitectónico central abordado y denunciado técnicamente por la presente tesis académica y logística documental. Al carecer las ONG cíclicas ciegamente y rígidamente de sólidas infraestructuras técnicas de información de núcleo empresarial \textit{Cloud-Native} y económicas, estas entidades filantrópicas se ven forzadas histórica e ineludiblemente a intentar burda coordinar y auditar cíclicamente todas y cada una de sus masivas e interconectadas transacciones operativas mediante el uso absolutamente inapropiado, inseguro y técnicamente contraproducente ciego de redes y sociales puramente comerciales y estáticamente domésticas, como lo son fundamental e indiscutiblemente los abigarrados y caóticos grupos ciegos públicos de \textit{Facebook} y los veloces pero terriblemente volátiles y ciegos e informales y muy y foros de difusión masiva en hilos tontos y estáticos de texto plano y plano documental de la app logística comercial popular de \textit{WhatsApp}.

Esta arcaica e improvisada solución de ruteo estática e indocumentada detona en el infame y severamente destructivo y fenómeno de fragmentación cíclico técnico conocido mundial e imperativamente en el rigor de la ingeniería global de datos y científicamente como el \textit{Data Siloing} e. Un y ciego estático "silo de información" cíclico ocurre inexorablemente y de manera letal cuando la data cíclica y sensible, invaluable e indispensable (tales como ciegos y métricos padrones médicos con números cíclicos vitales y o delicados montos y balances transaccionales públicos de abnegadas y altruistas donaciones) se encuentra estática y encerrada ciega e intermitentemente, amurallada y tecnológicamente en múltiples ecosistemas puramente comerciales e hiper-aislados (celulares personales con grupos ciegos independientes de WhatsApp y listas planas locales cíclicas de Excel y de portátiles domésticos e independientes) que resulta técnica e informáticamente puramente ineludible y ciego rígida e imposible de cruzar, auditar, exportar o analizar métrica y centralizadamente por el nodo y cíclico rígido comando central unificado de la misma propia y abnegada y filantrópica rígida organización ONG.

\subsection{El Impacto Crítico en Privacidad y Tolerancia a Fallos}

El cíclico ciego e implacable y uso y y letal abuso síncrono de los rudimentarios hilos ciegos estáticos documentales de texto puro y plano de redes sociales en las operaciones cíclicas y críticas logísticas de rescate y despachos no solo sepulta y hunde la agilidad de conteo y logística de manera trágicamente irreversible y. De manera infinitamente más preocupante y de severo e imperdonable y crítico rigor arquitectónico forense documental, expone pública y la privacidad logística y cíclica e intransferible de los abnegados donantes ciegos y voluntarios a catastróficas y vulnerables brechas de ruteo seguridad cibernética, al carecer éstas domésticas aplicaciones de auténtico estricto, robusto y sólido ciego rígido aislamiento \textit{Multi-Tenant} y carecer y técnicamente y de inquebrantables esquemas de Política de Seguridad de Filas (RLS) y relacionales ciegas estables. 

En un rudimentario e entorno de WhatsApp logístico, cualquier voluntario efímero tiene absurda, crítica y públicamente un indocumentado e indiscriminado acceso y completo a la gigantesca e rígida y ciega base e histórica ciega telefónica y base de nombres de todos los miembros; y es total y matemáticamente e incapaz e de generar cíclicamente transacciones técnicas métricas ACID estables rígidas que garanticen logísticamente el ciego balance seguro entre de sumas restas rígidas puramente de cíclica el inventario donado. Por tanto, la impostergable y urgente necesidad nacional e ingenieril ciega y de pura construcción técnica de un clúster en nube (\textit{Cloud Serverless}) con asombroso dinamismo cíclico \textit{Virtual DOM} y de ruteo \textit{WebSocket} propuesto centralmente en esta propia tesis universitaria no es un simple y mundano capricho y o un lujo estético; es verdaderamente y una letal y crítica y puramente e indispensable y urgente operación y deber ingenieril y humanitario rígido y civil cíclico que busca inquebrantablemente digitalizar, unificar de manera y blindar criptográfica y técnicamente y el futuro logístico de las nobles organizaciones y solidarias filantrópicas del civil tercer sector nacional.
